\section{Introduction}
\label{sec:yba-introduction}

Consider one memory register used by two tasks. Each task would set the register to a different value. Moving it towards one target moves it away from the other. The learner can change the register, change the relative authority of the tasks, or buy a second register. These are different operations. No amount of care in taking the gradient of a fixed objective will make the third operation appear if the architecture permits only the first two.

The difficulty becomes sharper when the learner must decide for itself which operation is worth making. Its estimate of future benefit depends on the model it currently has. Its estimate of the cost depends on the memory, controller and action repertoire through which that model is implemented. A proposed revision can alter both. The evaluator is therefore not merely a number attached to the model by an observer. It is a mechanism within the loop, with a state, inputs and an update protocol of its own.

This does not make a supplied loss illegitimate. An engineering problem can specify a target externally and ask an algorithm to optimise it. Our question begins where that description stops: how should we describe a learner that regulates its own model revision, including the allocation of a shared representation? We need to distinguish three things that are easily conflated. A controller may be internally consistent, it may converge, and it may improve an external task. None of these properties alone supplies the other two.

We develop the argument from the register rather than introducing a universal objective. Two quadratic demands let us calculate the cost of shared allocation exactly. An explicit split option then identifies when purchasing capacity is cheaper than continuing to arbitrate. The result is a price boundary, not a claim that a larger model must generalise better. Every option is compared on the same ledger.

The next step is to give the learner actual actions. A model revision must be scored at the state it will reach, not at an ideal fit that the implementation never executes. For a partial least-squares move, the difference can reverse the decision to update. With those actions specified, a finite relative-entropy objective produces a Gibbs selector relative to a declared reference measure. The statistical mechanics is a variational construction, not an assumption that the learner contains a thermostat or samples all its alternatives.

We then put an evaluator on the opening register. Its model coordinate and authority weight adapt together under a simple, explicitly postulated controller. This gives a concrete self-consistent loss geometry: a loss family, a metric, a model step and an evaluator step whose joint stationary states can be studied. The construction admits a rigorous convergence region. Its usefulness is precisely that we can say what converges and why, without identifying a stable internal agreement with successful regulation of the world.

A finite learner also needs descriptions smaller than its complete history. We therefore ask what remains when update opportunities are aggregated. Weighted free-energy contraction gives an exact answer for finite hidden alternatives. For the event that at least one of several independent opportunities produces an update, the retained quantity is cumulative hazard. Keeping only a summed margin and one common threshold is a further approximation. Its error can be measured, and its response to repricing has a simple meaning: the expected number of fine events conditional on any event occurring.

These calculations give the phrase \emph{laws of learning} its intended scope. We propose three requirements for this class of embodied model-revision schemes. \emph{Endogenous evaluation} requires an explicit description of the evaluator when it is part of the learner. \emph{Embodied price} requires a common ledger for the actions actually performed. \emph{Representation covariance} requires a declared receiving representation and a test of which predictions or transitions it preserves. The requirements organise a model class. The finite theorems follow from the premises stated beside them, not from a claim that every biological learner obeys our controller.

The paper is self-contained. Section~\ref{sec:yba-setup} derives the register purchase. Section~\ref{sec:self-consistent-loss-geometry} constructs the action selector and coupled regulator. Section~\ref{sec:partition-covariance-reuse} moves the construction inside a larger controller. Section~\ref{sec:numerical-demonstrations} examines a synthetic tracking problem and retains the finding that consistent action pricing does not deliver a predictive victory. Section~\ref{sec:yba-examples} develops exact contraction, the OR example and dynamical receiving-map tests. We then identify the primary methods used and develop the consequences for further modelling. The appendices provide proofs, the complete finite protocol and the quantum extension.
